% Study sheet: Volume I, Chapter 1, Introduction.
% Numbers: \val{...} from introduction_studysheet_numbers.tex, written by introduction_studysheet_numbers.py.
\documentclass{studysheet}
\input{introduction_studysheet_numbers}
\sheetvolume{1}
\sheetchapter{1}{Introduction}
\sheetsources{Numbers from the mlsysim registry and the chapter's scenarios; questions marked assumed use new inputs.}

\begin{document}
% ===================================================================== FRONT: LEARN
\frontside
\begin{bigidea}
An ML system has two jobs at once: its \hl{learned behavior} must stay trustworthy, and the
\hl{machine} must deliver it within limits of time, energy and cost. Because the behavior is
learned from data, it can get worse even when no code changes.
\end{bigidea}

\begin{leftcol}
\sheetsec{Key ideas}
\begin{ideas}
\idea{Behavior comes from data.}{When the world drifts from the training data, accuracy can fall
  while tests and health checks still pass: \emph{silent degradation}.}
\idea{Diagnose by axis: Data, Algorithm, Machine (D·A·M).}{Fix the axis that limits the result,
  then measure again; the bottleneck moves.}
\idea{Time goes to moving data, computing and overhead.}{A faster part helps only as much as its share (Amdahl's law).}
\idea{Moving data often costs more than computing on it.}{\emph{Arithmetic intensity}, operations
  per byte moved, says which one sets the time.}
\idea{The bitter lesson.}{Methods that keep improving with more computation win, given data and a
  machine to feed them.}
\idea{Three efficiency levers:}{fewer samples, fewer operations per sample, less time per operation.
  Judge each by accuracy gained per cost.}
\end{ideas}

\vspace{11pt}
\sheetsec{Three axes, and what each pair decides}
{\centering
\begin{tikzpicture}[sheet]
  \node[accbox, minimum width=60pt] (D) at (125,62) {Data};
  \node[accbox, minimum width=60pt] (A) at (40,12) {Algorithm};
  \node[accbox, minimum width=60pt] (M) at (210,12) {Machine};
  \draw[dim] (D) -- (A) (D) -- (M) (A) -- (M);
  \node[note, anchor=east] at (76,40) {what to learn from};
  \node[note, anchor=west] at (174,40) {how data moves};
  \node[note, anchor=north] at (125,7) {how to execute};
\end{tikzpicture}\par}

\vspace{10pt}
\fcard{Return on compute}{\text{RoC} = \frac{\Delta\,\text{accuracy}}{\Delta\,\text{cost}}}%
  {Compare upgrades by points of accuracy per dollar, not by accuracy alone.}{}{}

\vspace{11pt}
\sheetsec{Numbers to know}
\begin{tiles}
\tile{$\sim$10$^{\val{growth}}\times$}{growth in training compute since AlexNet}\hfill
\tile{1/\val{algox}}{compute EfficientNet-B0 needed to match AlexNet}
\end{tiles}
\end{leftcol}%
\begin{rightcol}
\sheetsec{Formulas you will use}
\fcard{Iron law}%
  {\Time = \under{\frac{\text{samples}}{\text{task}}}{Data} \times
           \under{\frac{\text{FLOPs}}{\text{sample}}}{Algorithm} \times
           \under{\frac{\text{seconds}}{\text{FLOP}}}{Machine}}%
  {Each factor is one lever, and changing one can move another.}{}{}
\fcard{Iron law, as time spent}%
  {\Time = \frac{\DataVolume}{\Bandwidth} + \frac{\Operations}{\PeakRate\,\EfficiencyHw} + \Latency}%
  {Bytes moved $\div$ bandwidth, plus operations $\div$ achieved rate ($\PeakRate$ peak $\times$
   $\EfficiencyHw$ share achieved), plus fixed overhead $\Latency$.}%
  {With overlap, the larger of the first two replaces their sum. Seconds per FLOP $= \Time / \Operations$.}{}
\fcard{Arithmetic intensity}%
  {\ArithmeticIntensity = \frac{\Operations}{\DataVolume}}%
  {Operations per byte moved. Compare it with the machine's balance $\PeakRate\,\EfficiencyHw / \Bandwidth$:
   below it, memory sets the time; above it, compute does.}{}%
  {\begin{tikzpicture}[sheet]
     \draw[dim] (10,0) -- (240,0);
     \foreach \tk/\tl in {0.1/0.1,1/1,10/10,100/100,1000/1000 pJ}{
       \pgfmathsetmacro\x{10+(log10(\tk)+1)*57.5}
       \draw[dim] (\x,0) -- (\x,-2.5);
       \node[small, anchor=north] at (\x,-4) {\tl};}
     \pgfmathsetmacro\xa{10+(log10(\val{efp})+1)*57.5}
     \pgfmathsetmacro\xb{10+(log10(\val{edram})+1)*57.5}
     \filldraw[fill=accentsoft, draw=accent, line width=0.5pt] (\xa,0) circle (2.4pt);
     \filldraw[fill=accent, draw=accent] (\xb,0) circle (2.4pt);
     \node[note, anchor=south] at (\xa,5) {16-bit operation \val{efp} pJ};
     \node[note, anchor=south] at (\xb,5) {DRAM byte \val{edram} pJ};
   \end{tikzpicture}}
\fcard{Degradation equation}%
  {\AccuracyAfterDeployment(\DeploymentAge) \approx \InitialAccuracy - \FittedSensitivity \cdot \StatisticalDivergence}%
  {Accuracy falls with the divergence between today's inputs $\CurrentDistribution$ and training
   inputs $\TrainingDistribution$; $\FittedSensitivity$ is fitted from outcomes. Drift warns; labels measure.}{}%
  {\begin{tikzpicture}[sheet, every node/.style={nbox, text width=46pt, align=center, minimum height=20pt}]
     \node (a) at (28,0) {world changes};
     \node (b) at (92,0) {inputs drift};
     \node[hotbox, text width=46pt, align=center, minimum height=20pt] (c) at (156,0) {accuracy may drop};
     \node (d) at (220,0) {labels measure it};
     \draw[arr] (a) -- (b); \draw[arr] (b) -- (c); \draw[arr] (c) -- (d);
   \end{tikzpicture}}
\end{rightcol}

% ===================================================================== BACK: PRACTICE
\backside
\begin{answers}
\ans{Finish}{\val{pre} + \val{modnew} + \val{post} = \val{anew} ms; \val{atot} $\div$ \val{anew} =
  \val{ax}$\times$ faster (\val{apct}\% less time); the other \val{prepost} ms did not change.}
\ans{1}{\val{qm} + \val{qc} + \val{qo} = \val{qseq} ms; overlapped \val{qmax} + \val{qo} = \val{qovl} ms.
  $\ArithmeticIntensity \approx \val{qi}$ FLOP/byte, above the balance \val{qtfs} $\div$ \val{qtbs} = \val{qbal} FLOP/byte, so compute
  sets it; at \val{qx}$\times$ math, compute \val{qcfast} ms $<$ move \val{qm} ms, so memory sets it: \val{qmo} ms.
  Energy: moving \val{qem} J vs math \val{qec} J, \val{qer}$\times$ more for moving.}
\ans{2}{(a) samples per task, Data; (b) seconds per FLOP, Machine; (c) FLOPs per sample, Algorithm. (d) No. Better algorithms make each FLOP do more; methods that scale can still turn added compute into
  accuracy, but only while data and a machine can feed them.}
\ans{3}{A point is worth \$\val{worth}, so each \$\val{per} must buy at least \val{rocbe} points. A buys
  \val{roca} per \$\val{per} and pays; B buys \val{rocb} and does not.}
\ans{4}{No. Now: compare this week's inputs with the training inputs (a drift warning) and keep checking
  latency and memory. Next week: measure accuracy on labeled outcomes, overall and per slice.}
\end{answers}

\begin{leftcol}
\sheetsec{Worked example}
\begin{practicebox}\boxtitle{Where does a faster accelerator help?}
A pass spends \val{m} s moving data, \val{c} s computing, \val{o} s on overhead.
New hardware makes the math \val{s}$\times$ faster.\par\vspace{5pt}
{\centering
\begin{tikzpicture}[sheet]
  \pgfmathsetmacro\sc{150/\val{tot}}
  \pgfmathsetmacro\pA{58+\val{m}*\sc}  \pgfmathsetmacro\pB{\pA+\val{c}*\sc}  \pgfmathsetmacro\pC{\pB+\val{o}*\sc}
  \pgfmathsetmacro\qB{\pA+\val{cnew}*\sc}  \pgfmathsetmacro\qC{\qB+\val{o}*\sc}
  % today
  \node[note, anchor=west] at (0,30) {today};
  \draw[nbox] (58,23) rectangle (\pA,37);  \node[lab] at ({(58+\pA)/2},30) {move};
  \draw[accbox] (\pA,23) rectangle (\pB,37); \node[lab] at ({(\pA+\pB)/2},30) {compute};
  \draw[nbox] (\pB,23) rectangle (\pC,37);
  \node[lab, anchor=west] at ({\pC+5},30) {= \val{tot} s};
  % after
  \node[note, anchor=west] at (0,8) {\val{s}$\times$ faster math};
  \draw[nbox] (58,1) rectangle (\pA,15);   \node[lab] at ({(58+\pA)/2},8) {move};
  \draw[hotbox] (\pA,1) rectangle (\qB,15); \node[lab, text=white] at ({(\pA+\qB)/2},8) {compute};
  \draw[nbox] (\qB,1) rectangle (\qC,15);
  \node[lab, anchor=west] at ({\qC+5},8) {= \val{new} s};
\end{tikzpicture}\par}\vspace{4pt}
\begin{steps}
\item Today: \val{m} s + \val{c} s + \val{o} s = \res{\val{tot} s}.
\item After: \val{m} s + \val{c} s $\div$ \val{s} + \val{o} s = \res{\val{new} s}, not \val{tot} s $\div$ \val{s} = \val{naive} s.
\end{steps}
\takeaway{An upgrade speeds up only the share of time it touches.}
\end{practicebox}

\vspace{12pt}
\sheetsec{Your turn: finish it}
\begin{practicebox}\boxtitle{A faster model inside a request}
A request spends \val{pre} ms preparing input, \val{mod} ms in the model, \val{post} ms after.
The model becomes \val{sp}$\times$ faster.\par\vspace{5pt}
{\centering
\begin{tikzpicture}[sheet]
  \pgfmathsetmacro\sc{236/\val{atot}}
  \pgfmathsetmacro\xa{\val{pre}*\sc} \pgfmathsetmacro\xb{\xa+\val{mod}*\sc} \pgfmathsetmacro\xc{\xb+\val{post}*\sc}
  \draw[nbox] (0,0) rectangle (\xa,16);   \node[lab] at ({\xa/2},8) {prepare \val{pre} ms};
  \draw[hotbox] (\xa,0) rectangle (\xb,16); \node[lab, text=white] at ({(\xa+\xb)/2},8) {model \val{mod} ms};
  \draw[nbox] (\xb,0) rectangle (\xc,16); \node[lab] at ({(\xb+\xc)/2},8) {after \val{post} ms};
\end{tikzpicture}\par}\vspace{5pt}
\textbf{Given:} new model time = \val{mod} $\div$ \val{sp} = \val{modnew} ms.\par
\textbf{You:} find the new request time, how many times faster the request is, and say why it is not \val{sp}$\times$.
\end{practicebox}

\vspace{12pt}
\sheetsec{Common mistakes}
\mistake{``High accuracy on the test set means it is ready to deploy.''}%
  {\val{vgtest} validation images sample a vanishing part of ${\sim}10^{\val{vgdigits}}$ possible images, and say
   nothing about drift or cost.}{}
\end{leftcol}%
\begin{rightcol}
\mistake{``The tests pass and latency is fine, so the model is fine.''}%
  {Learned behavior drifts with no code change: losing \val{drop} points a month, \val{acc0}\% becomes
   \val{accfinal}\% in \val{months} months.}%
  {\begin{tikzpicture}[sheet]
     % x: months 0..6 -> 30..236; y: 79..86 % -> 0..48
     \pgfmathsetmacro\sx{206/\val{months}}
     \foreach \m in {0,2,4,6}{\draw[gridline] ({30+\m*\sx},0) -- ({30+\m*\sx},48);
       \node[small, anchor=north] at ({30+\m*\sx},-3) {\m{} mo};}
     \foreach \yv in {80,85}{\draw[gridline] (30,{(\yv-79)*48/7}) -- (236,{(\yv-79)*48/7});
       \node[small, anchor=east] at (26,{(\yv-79)*48/7}) {\yv\%};}
     \draw[accent, line width=1.1pt] (30,{(\val{acc0}-79)*48/7}) -- (236,{(\val{accfinal}-79)*48/7});
     \filldraw[fill=accent, draw=accent] (236,{(\val{accfinal}-79)*48/7}) circle (2pt);
     \node[lab, anchor=north east] at (232,{(\val{accfinal}-79)*48/7-3}) {\val{accfinal}\%};
     \node[note, anchor=south west] at (34,{(\val{acc0}-79)*48/7+2}) {no code changed};
   \end{tikzpicture}}

\vspace{3pt}
\sheetsec{Check yourself}
\begin{qs}
\item \assumed{} a step moves \val{qgb} GB at \val{qtbs} TB/s and does \val{qgf} GFLOP of 16-bit math at \val{qtfs} TFLOP/s
  (achieved rates), plus \val{qo} ms overhead. Time with no overlap and with full overlap? What sets it,
  before and after \val{qx}$\times$ faster math? Which costs more energy, the bytes or the math?
\item Which iron-law factor changes when you (a) train on half as many, better-chosen samples, (b) run the
  same model in 8-bit on hardware that runs 8-bit operations faster (same operation count), (c) switch to a smaller architecture?
  (d) Compute for selected training runs grew about 10$^{\val{growth}}\times$ since AlexNet, yet EfficientNet-B0 reached
  AlexNet's accuracy with 1/\val{algox} of its compute. Does that contradict the bitter lesson? When does more compute help?
\item \assumed{} upgrade A adds \val{ga} points for \$\val{ca}; B adds \val{gb} points for \$\val{cb}.
  A point of accuracy is worth \$\val{worth}. Which upgrades are worth buying?
\item A classifier meets its latency budget. This week's inputs have shifted, and their labels arrive next
  week. Can you certify its accuracy today? What do you measure now, and what when the labels arrive?
\end{qs}
\end{rightcol}
\end{document}
